\section{Keamanan Web: Validasi Input, Sanitasi, dan Praktik Aman}

Keamanan web dimulai dari prinsip fundamental: \textbf{jangan pernah percaya data dari client}. Setiap input --- dari form, URL, cookie, atau header HTTP --- harus divalidasi dan disanitasi di server sebelum diproses. \textbf{Validasi} memeriksa apakah data sesuai format yang diharapkan (wajib diisi, panjang karakter, format email, rentang angka). \textbf{Sanitasi} membersihkan data agar aman digunakan --- misalnya \code{htmlspecialchars()} untuk output ke HTML, \code{strip\_tags()} untuk menghapus tag HTML. Validasi menolak data buruk; sanitasi memperbaiki data agar aman \cite{phpRightWay}.

Praktik aman lainnya: simpan password dengan \code{password\_hash(PASSWORD\_DEFAULT)} dan verifikasi dengan \code{password\_verify()}; jangan simpan password plain text. Gunakan HTTPS untuk enkripsi transport. Batasi error message di produksi --- jangan tampilkan stack trace atau detail database. Simpan file upload di luar web root atau dengan nama acak. Terapkan prinsip \textit{least privilege} untuk user database dan file permission \cite{owaspXss}.

\begin{webcode}[language=PHP, caption={Validasi dan sanitasi input}]
<?php
function sanitize(string $input): string {
  return htmlspecialchars(trim($input), ENT_QUOTES, 'UTF-8');
}

$errors = [];
$nama  = sanitize($_POST['nama'] ?? '');
$email = trim($_POST['email'] ?? '');
$umur  = intval($_POST['umur'] ?? 0);

if (strlen($nama) < 3) $errors[] = "Nama minimal 3 karakter.";
if (!filter_var($email, FILTER_VALIDATE_EMAIL)) {
  $errors[] = "Email tidak valid.";
}
if ($umur < 17 || $umur > 100) $errors[] = "Umur harus 17-100.";

// Password hashing
$hash = password_hash($_POST['password'], PASSWORD_DEFAULT);
// Verifikasi: password_verify($input, $hash)
?>
\end{webcode}

